\chapter{Why Reactions Release or Absorb Heat}

\begin{kaichang}
In a winter classroom, someone always pulls out a hand warmer: open the wrapper, rub it, and minutes later it warms you for hours. No fire, electricity, or battery---where does the heat come from?

Hospitals and sports grounds have the opposite pouch: an instant cold pack. Squeeze until an inner bag pops, and within seconds the whole pack feels icy against a sprained ankle. No refrigerator or compressor---where does the cold come from?

One apparently creates heat, the other cold. Guess whether the warmth was already hidden inside and whether the cold is manufactured. Write your answers down; by the end, we will balance both pouches' energy accounts.
\end{kaichang}

\section{First Distinguish System from Surroundings}

Holding a warmer, you feel that it produces heat. More precisely, \textbf{energy flows from the pouch into your hand}. This seemingly fussy distinction underpins the chapter.

Chemists call the small part being studied the \textbf{system}: iron powder and air inside a warmer, or reacting solution in a beaker. Everything else is the \textbf{surroundings}: your hand, the beaker, room air. If a reaction lowers system energy, excess energy flows outward as heat and warms the surroundings: an \textbf{exothermic} reaction. If system energy rises, it draws from the surroundings, cooling them: an \textbf{endothermic} reaction or process.

Thus saying a cold pack creates cold reverses the story. Your hand and the pouch surface cool because the process inside takes their heat. The pack \textbf{absorbs heat}, not manufactures cold.

Examples surround you:

\begin{itemize}[itemsep=2pt]
\item Methane combustion on a gas stove releases heat vigorously to boil water.
\item Quicklime, \chem{CaO}, releases heat with water, bringing self-heating hotpot broth to the boil.
\item Acid--base neutralization releases heat; dilute hydrochloric acid added to sodium hydroxide visibly warms the beaker.
\item Baking soda with vinegar absorbs heat, cooling the cup, as you will soon try.
\item Ammonium nitrate dissolving strongly absorbs heat: the cold pack's secret.
\end{itemize}

The last is not even a typical chemical reaction, merely \textbf{dissolving}. Heat flow therefore does not depend on whether a new substance forms. Whenever particle connections change, energy accounts may differ. Chapter 25 described dissolution as rearrangement; here we add price tags.

Dissolution heat works both ways. Ammonium nitrate absorbs, but sodium hydroxide releases considerable heat. Kitchen drain cleaner uses it to attack grease, and adding water makes the container hot. Concentrated sulfuric acid dilution is more vigorous still, able to splash and burn. Only teachers should demonstrate it; never handle it at home, and always add acid to water, never the reverse. The same act of adding something to water may cool, heat strongly, or do almost nothing. Two competing payments explain the difference, as we will see with cold packs.

Another dramatic example mixes barium hydroxide crystals and ammonium chloride in a beaker. It becomes cold enough to \textbf{freeze} a wet wooden board beneath it, so lifting the beaker lifts the board. Two stirred solids make ice without refrigeration. Toxic barium salts restrict this to teacher demonstrations or videos, but it vividly demonstrates endothermic reaction: heat taken from surroundings freezes water.

\section{Zooming In on Chemical Bonds}

Where do the price tags go? On bonds.

Recall Chapters 17--19: bonds are not sticks but \textbf{lower-energy arrangements} of joined atoms. Separating bonded atoms therefore requires \textbf{energy input}; joining them \textbf{releases energy}.

Use magnets as an analogy. Pulling attached magnets apart takes work, storing energy in their separated state. Release them and they snap together, jolting your hand as energy leaves. Chemical bonding substitutes electrical atomic/electronic interactions for magnetism:

\begin{itemize}[itemsep=2pt]
\item \textbf{Breaking absorbs}: breaking a bond always costs energy, never provides a free gain.
\item \textbf{Forming releases}: forming a bond always releases energy, never demands payment.
\end{itemize}

Chapter 26 sees atoms rearranging; energy accounting sees two transactions: pay to break old bonds, receive for forming new ones. Heat release or absorption depends on \textbf{the difference}:

\[ \text{Net energy change} \approx \text{Energy absorbed breaking bonds} - \text{Energy released forming bonds} \]

A negative result means receipts exceed expenses, system energy falls, and heat leaves. Positive means system energy rises and heat enters. Exothermic reactions do not release stored heat; \textbf{product bonds are stronger and more favorable than reactant bonds}, leaving the rearranged system at lower energy.

\section{A Real Calculation}

These prices are measurable. Chemists determine \textbf{bond energies}: average kilojoules needed to break one mole (\ud{6.02\times 10^{23}}{} bonds) in the gas phase. Some familiar examples:

\begin{table}[h]
\centering
\begin{tabular}{cc}
\toprule
Bond & Approximate energy (\ud{}{kJ\cdot mol^{-1}}) \\
\midrule
\chem{H-H} & 436 \\
\chem{O=O} & 498 \\
\chem{H-O} & 463 \\
\chem{C-H} & 413 \\
\bottomrule
\end{tabular}
\end{table}

Use Chapter 26's balanced hydrogen combustion:

\[ \chem{2H_2} + \chem{O_2} \rightarrow \chem{2H_2O} \]

Breaking 2 \chem{H-H} and 1 \chem{O=O} costs $2\times 436+498=1370$ kilojoules. Forming 4 \chem{H-O} bonds in two waters returns $4\times 463=1852$. Net: $1370-1852=-482$ kilojoules, about 482 kilojoules released per two moles of hydrogen.

That can heat roughly 1.4 kilograms of water from room temperature to boiling. Yet two moles of hydrogen weigh only four grams: a small handful of gas holds enough for a pot of soup. This is why rockets carry liquid hydrogen and oxygen: hydrogen ranks among the best fuels by heat released per weight.

\begin{qiaoliang}
Most reactions occur not in sealed steel vessels but in open beakers, pots, and air, at nearly constant atmospheric pressure. One complication: heat release or gas formation can expand the system, doing work to push air aside. Energy leaves as volume work, not the heat measured. Chemists pragmatically define \textbf{enthalpy}, $H$, combining internal energy with pressure times volume. Under \textbf{constant pressure with no other work}, exchanged heat equals enthalpy change:
\[ q_p = \Delta H \]
$\Delta H$, enthalpy change, is total product enthalpy minus total reactant enthalpy. Accounts use the \textbf{system's viewpoint}: $\Delta H<0$ means decreasing enthalpy and energy flowing out, \textbf{exothermic}; $\Delta H>0$ means increasing enthalpy and heat drawn in, \textbf{endothermic}. You need not calculate absolute $H$. Chemists care only about \textbf{changes}, like altitude differences with an arbitrary zero.
\end{qiaoliang}

Add a new column to Chapter 26's ledger: an equation with enthalpy becomes a \textbf{thermochemical equation}:

\[ \chem{2H_2}(g) + \chem{O_2}(g) \rightarrow \chem{2H_2O}(g) \qquad \Delta H = -484\ \ud{}{kJ} \]

The negative sign means heat release. The $g$ state symbols are essential: making liquid water and vapor releases different heat, since evaporation itself absorbs energy (Chapter 22's intermolecular account).

Picture energy vertically, reactants higher and products lower, separated by $\Delta H$. Exothermic reactions resemble downhill water; endothermic products sit higher, with surroundings paying the difference.

\begin{shiyan}
\textbf{Make a ``Cold Pack'': Baking Soda and Vinegar}

\textbf{Materials}: two spoonfuls of kitchen baking soda, half a cup of vinegar, two disposable cups, a thermometer (or clean fingers), and a chopstick.

\textbf{Procedure}: (1) Pour vinegar into a cup and measure or feel its temperature. (2) Add both spoonfuls of soda and stir gently. (3) Watch temperature changes and rest the back of your hand against the cup.

\textbf{What to observe}: while bubbles form, does temperature rise or fall? Does the cup feel warm or cool? After gas escape and reaction cease, what happens if the cup stands undisturbed?

This reaction, bicarbonate with acid producing carbon dioxide, water, and acetate, is overall \textbf{endothermic}, taking heat from vinegar, cup, and hand. Afterward it slowly warms as surroundings return heat, ultimately reaching the same temperature.

\textbf{Safety note}: these edible ingredients are very safe, but vinegar stings eyes, so stir carefully. \textbf{Never} shake the bubbling mixture in a sealed bottle: carbon dioxide could turn it into a small bomb. Rinse the finished liquid down the sink.
\end{shiyan}

\section{Energy Accounts Depend Only on Start and Finish}

Now for the chapter's most convenient and profound rule.

Travel from a town at a mountain's foot to its summit by winding road, steep climb, or detour through a valley. Whatever the route or meals along it, your \textbf{net altitude gain} depends only on the endpoints.

Enthalpy works that way: a \textbf{state function}, whose change depends on initial and final states, not the path. In 1840, the Russian chemist \textbf{Germain Hess} generalized calorimetry results into \textbf{Hess's law}: \textbf{a reaction's total enthalpy change is the same whether completed in one step or several}.

Its power appears immediately when heat cannot be measured directly. Carbon's incomplete combustion to carbon monoxide is hard to isolate cleanly without byproducts. Take a detour:

\begin{itemize}[itemsep=2pt]
\item Route A, one step: complete carbon combustion, $\chem{C} + \chem{O_2} \rightarrow \chem{CO_2}$, measured $\Delta H=-394$ kilojoules.
\item Route B, two steps: first $\chem{C} + \frac{1}{2}\chem{O_2} \rightarrow \chem{CO}$, with uncertain value $x$; then $\chem{CO} + \frac{1}{2}\chem{O_2} \rightarrow \chem{CO_2}$, measured at $-283$ kilojoules.
\end{itemize}

Identical endpoints require $x+(-283)=-394$, giving $x=-111$ kilojoules. Two measurable rulers assemble an unmeasurable heat. Carbon becoming diamond, glucose fermentation, and many slow or dangerous reactions have enthalpies calculated by detour without performing the actual reaction.

These paper detours invisibly underpin industry. Ammonia and sulfuric-acid plants face \textbf{heat management}: unreleased heat causes overheating and runaway; insufficient input stalls reactions. Engineers divide processes into steps, consult handbook and calorimetric enthalpies, add and subtract them using Hess's law, and size heat exchangers and cooling accordingly. The energy account balances on drawings before any accident. Chapters 30 and 31 explain the further care with temperatures and catalysts.

\begin{zhengju}
Hess's law may seem obvious from energy conservation, but in 1840 it was not. Caloric theory still dominated, treating heat as a weightless fluid called caloric: more meant hotter, less colder. Heat was a \textbf{substance}, and releasing it had no apparent reason to be path-independent.

Hess relied on experiments with his improved calorimeter, measuring neutralization and dissolution. Total heat matched however many steps a route took. He elevated this from observation to a universal rule, strange under caloric theory but later a beautiful consequence of energy conservation: heat is not matter but \textbf{energy in transfer}, and energy transforms without appearing or vanishing.

Caloric theory's undoing also had a vivid experiment. In 1798, Count Rumford supervised cannon boring and noticed seemingly inexhaustible frictional heating. If caloric were stored inside metal, surely removed shavings should eventually exhaust it? Yet little material was removed, and even a blunt drill producing almost no shavings released plentiful heat. Work could generate heat \textbf{continuously}; it was energy flow, not stocked goods.

Measuring tools evolved too. In the 1780s, Lavoisier and Laplace built an ice calorimeter, placing a reaction---or even a guinea pig---inside an ice-surrounded chamber and weighing melted ice to measure heat. It showed that guinea-pig heat and oxygen consumption followed the same accounts as charcoal combustion: \textbf{life's warmth and furnace heat are the same thing}. This first bridge from chemistry to biology carries us through the book's latter half.
\end{zhengju}
\section{Measuring Reaction Heat: Calorimetry}

Hess's data came from \textbf{calorimeters}, not thin air. The principle is charmingly simple: react in a little room that retains heat, and measure its temperature change.

A school version uses an insulated or foam cup, thermometer, and stirring rod. Carry out, say, neutralization, measure warming, and calculate. Raising one gram of water by $1\,^{\circ}\mathrm{C}$ takes roughly 4.2 joules, its specific heat capacity:

\[ q = m \times 4.2\ \ud{}{J/(g\cdot ^{\circ}C)} \times \Delta T \]

One hundred grams of solution warming $5.4\,^{\circ}\mathrm{C}$ absorbs about $100\times4.2\times5.4\approx2270$ joules from the reaction. Divide by reactant moles for molar enthalpy. Dilute hydrochloric acid neutralizing sodium hydroxide gives about $-57$ kilojoules per mole of water formed.

Food needs a sturdier \textbf{bomb calorimeter}. Dried food goes into a pressure-resistant sealed steel bomb filled with oxygen, then burns completely. Surrounding water's temperature rise measures its chemical energy. The ``2000 kilojoules'' on biscuit packets comes from such burning. Your body instead uses dozens of gentle enzyme-led steps, discussed in Volume Three. But Hess's law gives exactly the same total for identical endpoints: biscuit plus oxygen to carbon dioxide plus water. Your body is a bomb calorimeter taking the winding road.

\section{The Model's Boundaries}

Enthalpy is powerful, but several things remain outside it:

\begin{itemize}[itemsep=2pt]
\item \textbf{$\Delta H$ is not reaction direction.} Exothermic means an energy surplus, yet cold-pack ammonium nitrate spontaneously dissolves while absorbing heat. Nature also keeps a dispersal account: Chapter 28's entropy.
\item \textbf{$\Delta H$ is not reaction speed.} Hydrogen and oxygen could release 484 kilojoules yet coexist for years until a spark triggers them. The energy drop determines whether it is worthwhile; the starting barrier whether it can get moving. Chapter 29 handles barriers.
\item \textbf{Enthalpy values have conditions.} Constant pressure is assumed, and temperature, pressure, and phase change values. Handbooks usually give standard changes at $25\,^{\circ}\mathrm{C}$ and one atmosphere with substances in their most stable states. Other settings need corrections.
\item \textbf{Bond energies are average estimates.} A table's 463 kilojoules for \chem{H-O} averages unequal costs for breaking water's first and second \chem{H-O} bonds. Errors of several percent are normal; precise values require calorimetry.
\end{itemize}

\begin{wuqu}
\textbf{``Energy is stored in chemical bonds and released by breaking them.''} Perhaps chemistry's most widespread falsehood, repeated even in popular-science books.

The table refutes it: breaking a mole of \chem{H-H} requires \textbf{input} of all 436 kilojoules. Otherwise hydrogen would split itself and generate electricity, leaving a very different universe.

Two analogies restore the direction. Magnets require work to separate and release energy when joining; energy resides in separation, not connection. Saying bonds store energy is like saying valley bottoms store height: energy belongs at the higher level. Hydrogen burns not because breaking releases energy, but because \chem{H-O} bonds are stronger than \chem{H-H} and \chem{O=O}, returning more from formation than breaking costs. \textbf{The net difference releases heat, never bond breaking itself.} The misconception returns disguised as ATP high-energy bond claims in Volume Three, where we dismantle it again.
\end{wuqu}

\section{Back to the Two Pouches}

Now fulfill our opening promise.

\textbf{Hand warmers} contain iron powder, activated carbon, salt, water, and insulating vermiculite. Open the wrapper, air enters, and iron rusts:

\[ \chem{4Fe} + \chem{3O_2} \rightarrow \chem{2Fe_2O_3} \qquad \Delta H \approx -1650\ \ud{}{kJ} \]

Four moles, about 224 grams of iron, release around 1650 kilojoules. A warmer's 30--40 grams can release over two hundred kilojoules, enough to heat two liters almost to boiling. It merely feels warm because salt and water tune the rate and vermiculite retains and slowly releases the heat over hours. Warmth was not hidden inside; it comes from \textbf{stronger iron--oxygen bonds}. Reactants existed beforehand but the wrapper kept them apart.

\textbf{Cold packs} contain ammonium nitrate and a small water sachet. Breaking it dissolves \chem{NH_4NO_3}. Pulling ions from the crystal costs energy, breaking Chapter 18's ionic bonding; surrounding ions with water returns energy through Chapter 25's solvation. For most salts, such as table salt, these nearly balance and temperature barely changes. Ammonium nitrate is exceptional: lattice separation costs more than hydration returns, net about $+26$ kilojoules per mole. The difference comes from solution, pouch, and sprained ankle. Your ankle cools because heat fills the dissolution deficit.

Hot or cold, one ledger: the difference between breaking old and making new connections determines heat flow, without exception.

Both pouches are legitimate medical and sports products. Cold constricts local vessels and slows swelling; warmth promotes circulation and relieves muscular stiffness. These principles have medical support, but heed the book's repeated rule: duration, temperature, and suitability require product instructions and medical advice. Frostbite and low-temperature burns are real risks. Chemistry supplies heat flow; humans are responsible for use.

\section{Another Account Beyond Energy}

\textbf{Translator safety note:} Do not put a hand inside a cold pack, open one for its chemicals, or perform the ammonium-nitrate exercise at home. Use intact packs only as labeled; leaking contents can be harmful. Treat the exercise as a paper design. See \href{https://www.poison.org/articles/whats-inside-ice-packs-201}{Poison Control's ice-pack guidance}.

You may have spotted a gap: if endothermic means rising system energy, why does ammonium nitrate readily dissolve? Put your hand in the cold pack---does dissolution stop? No, it continues to saturation, as though pushed toward \textbf{higher} energy.

This crucial part lies beyond our model: falling energy is not nature's only preference. Particles and energy favor spreading out; ordered ions often disperse into water even at an endothermic cost. This dispersal ledger is entropy. Together with enthalpy, it determines \textbf{which direction a process takes}. Next chapter calculates it.

\begin{tiaozhan}
\textbf{How Different Are Enthalpy and Internal Energy?} (Optional.) Internal energy, $U$, counts particle kinetic and interaction energy without pushing-air accounts. At constant pressure, $\Delta H=\Delta U+\Delta(pV)$. For ideal-gas reactions, $\Delta H=\Delta U+\Delta n_{\text{gas}}RT$, where $\Delta n_{\text{gas}}$ is the change in gas moles, $R$ the gas constant, and $T$ temperature. Hydrogen combustion to gaseous water gives $\Delta n_{\text{gas}}=2-3=-1$. At $25\,^{\circ}\mathrm{C}$, $\Delta n_{\text{gas}}RT\approx-2.5$ kilojoules, only about 0.5\% difference between enthalpy and internal-energy changes. Reactions neither producing nor consuming gas, such as solution neutralization, have nearly equal values. Chemists favor enthalpy because it quietly handles pushing the atmosphere aside.
\end{tiaozhan}

\section*{Try It}

\begin{sikaoti}
\item Draw energy diagrams for methane combustion (exothermic) and ammonium nitrate dissolution (endothermic). Mark reactants, products, and signs of $\Delta H$. Note that these show energy relationships, not actual particle photographs.
\item Estimate the energy to separate one mole of methane, \chem{CH_4}, with four \chem{C-H} bonds into gaseous atoms. Look up a measured value and consider why it differs.
\item A cup neutralization mixes 100 grams dilute hydrochloric acid and 100 grams sodium hydroxide solution, warming from $25.0\,^{\circ}\mathrm{C}$ to $28.2\,^{\circ}\mathrm{C}$. Estimate released heat using water's 4.2 \ud{}{J/(g\cdot ^{\circ}C)} specific heat. If the same reaction used 200 grams of solution, how would warming change?
\item A used warmer becomes hard and heavier. Explain the extra mass in particle terms and why conservation holds.
\item Is ``an exothermic reaction needs no heating'' correct? Give rebuttals from this chapter and the hydrogen--oxygen mixture, distinguishing exothermic, spontaneous, and fast.
\item Design a comparison of temperatures when equal amounts of salt and ammonium nitrate dissolve. State controlled variables, predictions, and the two competing payments behind their difference.
\end{sikaoti}
